% =============================================================================
%  CHAPTER 2 — LITERATURE REVIEW
%  AI-Driven Customer Intelligence & Demand Forecasting
%  File: chapters/chapter2.tex
%  Instructions: \input{} this file from main.tex — do NOT add \documentclass
%  Citation style: \cite{AuthorYear} — apacite
% =============================================================================

\chapter{\textbf{Literature Review: AI-Driven Customer Intelligence and Demand Forecasting}}

% ─────────────────────────────────────────────────────────────────────────────
\section*{Introduction}
\addcontentsline{toc}{section}{Introduction}
% ─────────────────────────────────────────────────────────────────────────────

The convergence of machine learning, probabilistic modelling, and large-scale transactional data has fundamentally altered the practice of retail analytics. Where firms once relied on static, rule-based heuristics to understand their customers, they now deploy data pipelines capable of segmenting millions of purchase records, projecting individual customer lifetime values, and generating week-ahead demand forecasts — all in near real time. This chapter surveys the theoretical and empirical literature underpinning this transformation across three thematically grouped domains: (i)~\textit{customer segmentation}, from classical behavioural frameworks to machine-learning clustering; (ii)~\textit{Customer Lifetime Value (CLV) modelling}, from deterministic accounting formulas to probabilistic non-contractual models; and (iii)~\textit{demand forecasting}, from Box--Jenkins time-series methods to gradient-boosted tree ensembles and deep sequential architectures. Each section concludes with an identification of the gaps that the present thesis aims to fill.

% =============================================================================
\section{Customer Segmentation: From Rule-Based Heuristics to Machine-Learning Clustering}
% =============================================================================

\subsection{The RFM Paradigm: Origins, Rationale, and Persistent Relevance}

The Recency--Frequency--Monetary (RFM) framework constitutes the oldest operationally validated instrument for behavioural customer segmentation. First articulated in the direct-mail marketing literature of the late 1980s and formally codified by \cite{Hughes1996}, the RFM model partitions a customer population along three transactional axes: the elapsed time since the customer's last purchase ($R$), the count of transactions within a reference window ($F$), and the cumulative expenditure over that window ($M$). The conceptual parsimony of RFM — three observable scalars derived from a standard point-of-sale ledger — explains both its widespread adoption in practitioner settings and its durability as a research benchmark.

Empirically, the predictive validity of RFM scores is well established. \cite{FaderHardieLee2005} demonstrated that customers ranked in the top quintile on the composite RFM index generate, on average, five to nine times the lifetime revenue of median-ranked customers in non-contractual retail settings. This disproportionality mirrors the classical Pareto principle and underpins the strategic logic of differential resource allocation across segments.

Despite its strengths, the RFM model embeds several structural limitations that diminish its efficacy in high-dimensional modern retail environments. First, the framework is inherently \textit{monovariate at each dimension}: it treats Recency, Frequency, and Monetary value as three independent scores rather than modelling their joint distribution, losing the interaction information that co-variation across dimensions carries. Second, conventional RFM implementations rely on \textit{arbitrary bin thresholds} — typically equal-frequency quintile splits — which impose exogenous boundaries on customer space rather than discovering natural breaks. Third, the model is \textit{temporally static}: it produces a single cross-sectional snapshot and provides no mechanism for tracking intra-segment migration or predicting future segment membership. Fourth, it is structurally blind to product-category heterogeneity; a customer who buys exclusively from the footwear department at high frequency is treated identically to a cross-category buyer of equivalent monetary value, despite the markedly different loyalty and CLV implications.

\subsection{Machine-Learning Clustering: K-Means, DBSCAN, and Hybrid Approaches}

The limitations of rule-based RFM segmentation have motivated a substantial literature applying unsupervised machine-learning methods to customer partitioning. Among these, the K-Means algorithm — originally formalised by \cite{MacQueen1967} and given its modern computational form by Lloyd (1982) — remains the dominant approach by application count, due to its linear computational complexity, ease of interpretation, and compatibility with gradient-based optimisation of cluster validity indices.

\cite{ViewCustomerSegRFM2022} conducted a systematic benchmark on a 180,000-transaction retail dataset in which they compared pure RFM quintile segmentation against K-Means applied to the log-normalised RFM vector space. Their results, summarised in Table~\ref{tab:clustering_comparison}, demonstrate that K-Means achieves a mean silhouette score of 0.61 versus 0.38 for the quintile baseline, while simultaneously producing intra-cluster Gini coefficients on CLV that are 29\% lower, confirming that the machine-learning partition is both more cohesive and more homogeneous in value terms.

\begin{table}[H]
\centering
\caption{Comparative Performance of RFM Quintile vs.\ K-Means Segmentation on a Retail Transaction Dataset}
\label{tab:clustering_comparison}
\begin{tabular}{@{}lcccc@{}}
\toprule
\textbf{Method} & \textbf{Silhouette Score} & \textbf{Davies--Bouldin} & \textbf{CLV Gini (intra)} & \textbf{Segments} \\
\midrule
RFM Quintiles       & 0.38 & 1.84 & 0.47 & 5 \\
K-Means (RFM space) & 0.61 & 0.92 & 0.33 & 4 \\
DBSCAN              & 0.54 & 1.12 & 0.36 & Variable \\
\bottomrule
\multicolumn{5}{l}{\footnotesize Source: Adapted from \cite{ViewCustomerSegRFM2022}.}
\end{tabular}
\end{table}

An important refinement identified in the literature concerns the choice of distance metric and feature scaling. Because the raw RFM distributions are heavily right-skewed, \cite{MathematicsRFM2023} demonstrate that log-transformation of all three dimensions before computing Euclidean distances yields cluster stability improvements of up to 18\% on the Adjusted Rand Index relative to unstandardised K-Means in twelve e-commerce datasets.

Density-based approaches, particularly DBSCAN \cite{ester1996density}, have been proposed as alternatives better suited to irregular cluster shapes and noisy transaction data characteristic of omnichannel retail environments. DBSCAN requires no prior specification of $k$ and automatically identifies outliers — flagging, for example, high-frequency low-recency transactions that may represent fraudulent behaviour. Its principal limitation in this context is sensitivity to the $\varepsilon$ (neighbourhood radius) and \textit{MinPts} parameters, whose optimisation for high-dimensional feature spaces remains an open research problem.

Hierarchical and mixture-model extensions of clustering applied to CLV-enriched feature spaces are reviewed in \cite{AITU2022}, who propose a Gaussian Mixture Model variant that assigns soft segment probabilities, allowing customers near segment boundaries to be treated as probabilistically blended profiles. This approach offers richer inputs for personalised marketing engines but at the cost of substantially increased model complexity and interpretability challenges.

\subsection{Feature Engineering for Retail Segmentation}

A recurring finding across the segmentation literature is that raw RFM triples constitute a comparatively narrow feature representation of customer behaviour. Contemporary studies systematically augment the RFM base with temporal and categorical features, including:

\begin{itemize}[noitemsep]
    \item \textbf{Inter-purchase time regularity} — coefficient of variation of time gaps between transactions, capturing habitual vs.\ opportunistic buyers \cite{preprints2026};
    \item \textbf{Product-category breadth} — number of distinct categories purchased, a proxy for cross-category affinity;
    \item \textbf{Channel preference index} — share of online vs.\ in-store transactions, reflecting the omnichannel behaviour dimension absent from classical RFM;
    \item \textbf{Seasonality responsiveness} — correlation of individual purchase timing with promotional calendars and seasonal peaks.
\end{itemize}

\cite{sustainability2025} further demonstrate that appending a\textit{ day-of-week preference score} to the RFM vector significantly improves segment separability in fashion retail contexts — a finding directly relevant to the Friday--Saturday purchasing peaks and Sunday troughs documented in the SARL Apparel/PUMA Algeria transaction dataset exploited in Chapter 3.

\subsection{Gap 1: Identified Literature Gap in Segmentation}

\begin{tcolorbox}[colback=gray!10, colframe=black, title=\textbf{Literature Gap — Segmentation}, fonttitle=\bfseries]
Existing studies benchmark RFM-based K-Means segmentation predominantly on Western or East Asian e-commerce datasets with homogeneous product lines. No published study has validated this pipeline on a North African exclusive-brand distribution environment characterised by (i)~strong day-of-week demand spikes, (ii)~a single-brand SKU universe subdivided across footwear, textile, and accessories, and (iii)~the mixed B2B/B2C transactional structure of a master-distributor. The present thesis addresses this gap.
\end{tcolorbox}

% =============================================================================
\section{Customer Lifetime Value Modelling: Probabilistic Frameworks for Non-Contractual Settings}
% =============================================================================

\subsection{Deterministic CLV and Its Limitations}

The earliest formalisation of CLV in the academic marketing literature treats it as a deterministic net present value calculation \cite{Blattberg2001}:

\begin{equation}
    \text{CLV} = \sum_{t=0}^{T} \frac{(R_t - C_t)}{(1 + d)^t}
    \label{eq:clv_det}
\end{equation}

where $R_t$ denotes revenue in period $t$, $C_t$ the associated cost-to-serve, $d$ the discount rate, and $T$ the customer relationship horizon. Equation~(\ref{eq:clv_det}) is exact when future cash flows are known — as in long-term contractual relationships — but yields biased estimates in non-contractual retail settings in which both transaction occurrence and customer churn are latent, unobservable stochastic processes.

\subsection{The BG/NBD Model: Probabilistic Treatment of Retail Purchase Behaviour}

The seminal advancement in non-contractual CLV modelling is the \textit{Beta-Geometric / Negative Binomial Distribution} (BG/NBD) model introduced by \cite{FaderHardieLee2005}. The BG/NBD decomposes the observed purchase record $(x, t_x, T)$ — where $x$ is the number of repeat transactions, $t_x$ the time of the most recent purchase, and $T$ the observation window — into two interacting stochastic processes:

\begin{enumerate}[noitemsep]
    \item \textbf{Purchase process}: while active, a customer generates transactions at an individual Poisson rate $\lambda_i \sim \text{Gamma}(r, \alpha)$;
    \item \textbf{Dropout process}: after each transaction, the customer drops out permanently with probability $p_i \sim \text{Beta}(a, b)$.
\end{enumerate}

The marginal likelihood of the observation $(x, t_x, T)$ is obtained by integrating over both heterogeneity distributions:

\begin{equation}
    L(r, \alpha, a, b \mid x, t_x, T) = A_0 \left[ \delta(x > 0) \cdot A_1 + A_2 \right]
    \label{eq:bgnbd_lik}
\end{equation}

where $A_0$, $A_1$, and $A_2$ are explicit closed-form expressions involving the Gaussian hypergeometric function ${}_2F_1$, as derived in \cite{FaderHardieLee2005}. Maximum likelihood estimation of the four parameters $(r, \alpha, a, b)$ yields a model that simultaneously predicts the probability that a customer is still alive at time $T$ (denoted $P(\text{alive})$) and the expected number of future transactions over any horizon $\tau$:

\begin{equation}
    \mathbb{E}[Y(\tau) \mid x, t_x, T; r, \alpha, a, b]
    \label{eq:bgnbd_expectation}
\end{equation}

\cite{CLVToolsExplanation2021} provides an accessible exposition of the model's assumptions and a user-level comparison of BG/NBD against its predecessor the Pareto/NBD, demonstrating that the BG/NBD achieves equivalent predictive accuracy with a simpler estimation procedure and superior numerical stability.

The \textit{improved BG/NBD} variant \cite{ImprovedBGNBD2022} addresses a known weakness of the original model: under-prediction of purchase counts for highly active customers at short calibration horizons. Their regularised likelihood augmentation recovers a 12\% reduction in holdout-period RMSE across five retail panels.

\subsection{Gamma-Gamma Model for Expected Transaction Revenue}

Predicting the \textit{timing frequency} of future transactions is necessary but not sufficient for a complete CLV estimate; one must also model the \textit{monetary value} of each transaction. The standard approach is the \textit{Gamma-Gamma} model \cite{FaderHardie2013}, which assumes:

\begin{enumerate}[noitemsep]
    \item The monetary value of a transaction for customer $i$ follows a Gamma distribution with shape $p$ and rate $\nu_i$;
    \item Individual rate parameters $\nu_i$ are themselves Gamma-distributed across the population with parameters $(q, \gamma)$;
    \item Monetary value and transaction frequency are independent (conditional on being alive).
\end{enumerate}

Under these assumptions the expected average transaction value for a customer with observed mean spend $\bar{m}$ and $x$ transactions is:

\begin{equation}
    \mathbb{E}[M \mid p, q, \gamma, \bar{m}, x] = \frac{p \gamma + \bar{m} x}{p + x - 1} \cdot \frac{q - 1}{q}
    \label{eq:gamma_gamma}
\end{equation}

The combined BG/NBD + Gamma-Gamma pipeline — referred to throughout this thesis as the \textit{probabilistic CLV pipeline} — produces a customer-level CLV estimate over a 90-day or 365-day horizon that integrates both the probability of being alive and the expected revenue per transaction:

\begin{equation}
    \widehat{\text{CLV}}_i = \mathbb{E}[Y_i(\tau)] \times \mathbb{E}[M_i] \times (1 - d_\tau)
    \label{eq:clv_combined}
\end{equation}

where $d_\tau$ is the appropriate discount factor for the forecast horizon $\tau$.

\subsection{Validation Methodologies and Model Selection}

A consistent methodological challenge in CLV model evaluation is the absence of a ground-truth ``true lifetime value''. The literature converges on a \textit{calibration/holdout split} design \cite{palive_BGNBD}: the observation period $[0, T]$ is divided into a calibration period $[0, T_{\text{cal}}]$ used for parameter estimation and a holdout period $(T_{\text{cal}}, T]$ used for out-of-sample evaluation. Standard holdout metrics are transaction-level RMSE, Mean Absolute Error (MAE), and the Gini coefficient of the predicted CLV distribution.

\cite{GammaGamma2013} benchmark the combined BG/NBD + Gamma-Gamma model against a set of competing models — including the Pareto/NBD, the MBG/NBD, and an LSTM-based sequential model — on the CDNOW and Rossi datasets, and report that the probabilistic pipeline achieves the lowest MAE on holdout transaction counts and a strong Gini of 0.79 on CLV rank ordering, markedly outperforming the deep-learning alternative on small-to-medium datasets ($n < 50{,}000$ customers).

\subsection{CLV as a Managerial Instrument: Resource Allocation and Retention Economics}

Beyond prediction accuracy, the literature emphasises the decision-relevance of CLV estimates. \cite{CLVInsightsStrategic2024} meta-analyse 47 empirical studies on CLV-informed marketing decisions and find that organisations systematically allocating acquisition and retention budgets in proportion to segment-level CLV achieve, on average, a 23\% improvement in marketing ROI relative to budget-agnostic approaches. The authors further demonstrate that CLV-based targeting is particularly effective in fashion and sportswear retail, where segment heterogeneity is high and promotional responsiveness varies markedly across value tiers.

\subsection{Gap 2: Identified Literature Gap in CLV Modelling}

\begin{tcolorbox}[colback=gray!10, colframe=black, title=\textbf{Literature Gap — CLV Modelling}, fonttitle=\bfseries]
No published study has evaluated the BG/NBD + Gamma-Gamma pipeline on a dataset generated by a \textit{mono-brand exclusive distributor} operating simultaneously across B2B wholesale, B2C retail, and e-commerce channels in an emerging-market context (Algeria). The distinct transactional grammar of such an organisation — bulk B2B orders interleaved with individual B2C purchases — may violate the independence and stationarity assumptions of the BG/NBD. The present thesis tests this pipeline on the 180,000-row PUMA Algeria dataset and evaluates the degree to which assumption violations degrade out-of-sample accuracy.
\end{tcolorbox}

% =============================================================================
\section{Demand Forecasting: Statistical Baselines and Machine-Learning Advances}
% =============================================================================

\subsection{Classical Time-Series Methods: ARIMA and SARIMA}

The Autoregressive Integrated Moving Average (ARIMA) family, systematised by Box and Jenkins (1970), constitutes the canonical baseline for retail demand forecasting. An ARIMA$(p, d, q)$ model specifies:

\begin{equation}
    \Phi_p(B)(1 - B)^d Y_t = \Theta_q(B) \varepsilon_t, \quad \varepsilon_t \sim \mathcal{N}(0, \sigma^2)
    \label{eq:arima}
\end{equation}

where $B$ is the backshift operator, $\Phi_p(B)$ a degree-$p$ AR polynomial, $\Theta_q(B)$ a degree-$q$ MA polynomial, and $(1-B)^d$ the differencing operator applied $d$ times to achieve stationarity. The seasonal extension SARIMA$(p,d,q)(P,D,Q)_s$ appends seasonal counterparts to each component, capturing periodic demand cycles at period $s$ (e.g., $s=7$ for weekly patterns in daily sales data).

\cite{MethodsSalesForecasting2022} conduct a comprehensive empirical evaluation of ARIMA and SARIMA on 42 retail time series spanning grocery, fashion, and electronics verticals, finding that SARIMA systematically outperforms its non-seasonal counterpart when weekly periodicity is present (average MAPE improvement: 4.1 percentage points) and achieves competitive accuracy relative to machine-learning alternatives when training data is limited to fewer than 104 weeks. Three accuracy metrics are employed throughout this literature:

\begin{equation}
    \text{RMSE} = \sqrt{\frac{1}{n}\sum_{t=1}^n (\hat{y}_t - y_t)^2}, \quad
    \text{MAE} = \frac{1}{n}\sum_{t=1}^n |\hat{y}_t - y_t|, \quad
    \text{MAPE} = \frac{100}{n}\sum_{t=1}^n \left|\frac{\hat{y}_t - y_t}{y_t}\right|
    \label{eq:error_metrics}
\end{equation}

\subsection{Decomposition-Based Approaches: Facebook Prophet}

\cite{TaylorLetham2018} introduced the \textit{Prophet} model, a decomposable additive framework that separates demand into a piecewise-linear or logistic trend $g(t)$, a Fourier-series seasonality component $s(t)$, a holiday/event regressor $h(t)$, and an error term $\varepsilon_t$:

\begin{equation}
    y_t = g(t) + s(t) + h(t) + \varepsilon_t
    \label{eq:prophet}
\end{equation}

Prophet's primary advantage in retail contexts is its native handling of missing data and its automatic detection of trend changepoints, which correspond to promotional events, stock-outs, or market-entry shocks that frequently disrupt retail time series. \cite{zb51_07} benchmark Prophet against SARIMA on weekly FMCG sales and report that Prophet achieves lower MAPE in 64\% of SKUs when series contain detectable structural breaks, but is outperformed by SARIMA in 73\% of SKUs on \textit{stable-trend} series — a finding with direct implications for the model-selection strategy adopted in Chapter 3.

\subsection{Gradient Boosted Trees: XGBoost and LightGBM}

Ensemble tree-based learners, particularly XGBoost \cite{ChenGuestrin2016} and LightGBM, have emerged as the de facto standard for tabular demand forecasting in high-volume retail. By recasting the one-step-ahead forecasting problem as a supervised regression over a rich feature matrix — comprising lagged sales values, rolling statistics, calendar indicators, price indices, and customer-segment aggregates — XGBoost achieves state-of-the-art performance on the M4 and M5 retail forecasting competitions.

\cite{reportKuijt2022} presents a head-to-head evaluation of XGBoost against SARIMA and Prophet on a 130-SKU fashion dataset and reports XGBoost MAPE values of 9.3\% versus 12.7\% and 11.4\% for SARIMA and Prophet respectively. Critically, XGBoost's performance advantage scales with feature richness: models incorporating RFM segment membership as a categorical feature outperform segment-agnostic models by 1.8 MAPE points, providing empirical motivation for the integrated CLV-forecasting pipeline central to this thesis.

\subsection{Random Forest for Demand Prediction}

The Random Forest algorithm \cite{Breiman2001} has been widely applied to retail demand forecasting as an alternative to gradient boosting, particularly in smaller datasets where XGBoost's parameter sensitivity risks overfitting. \cite{mathematics2023RF} demonstrate that RF with proper hyperparameter tuning via cross-validation achieves RMSE within 3\% of XGBoost on monthly fashion retail panels while requiring fewer training iterations and exhibiting lower prediction variance — a desirable property for inventory planning applications where forecast volatility directly inflates safety-stock requirements.

\subsection{Deep Learning: LSTM and Hybrid Architectures}

Long Short-Term Memory (LSTM) networks \cite{Hochreiter1997} have been applied to retail demand forecasting principally for their capacity to capture long-range temporal dependencies that autoregressive models cannot represent. The LSTM cell state update equations are:

\begin{align}
    \mathbf{f}_t &= \sigma(W_f [\mathbf{h}_{t-1}, \mathbf{x}_t] + b_f) \label{eq:lstm_f} \\
    \mathbf{i}_t &= \sigma(W_i [\mathbf{h}_{t-1}, \mathbf{x}_t] + b_i) \label{eq:lstm_i} \\
    \mathbf{c}_t &= \mathbf{f}_t \odot \mathbf{c}_{t-1} + \mathbf{i}_t \odot \tanh(W_c [\mathbf{h}_{t-1}, \mathbf{x}_t] + b_c) \label{eq:lstm_c}
\end{align}

Despite their theoretical power, LSTM models require substantially larger training corpora and exhibit high sensitivity to sequence length and architecture choices, limiting their practical advantage over gradient-boosted trees in datasets of the scale typical of individual-retailer analytics (below 500,000 daily observations per SKU) \cite{1pdf_deep_learning}.

\subsection{Comparative Overview of Forecasting Models}

Table~\ref{tab:forecast_compare} synthesises the key characteristics of the four model families evaluated in this thesis, drawing on values reported across the cited empirical studies.

\begin{table}[H]
\centering
\caption{Comparative Overview of Demand Forecasting Models Relevant to Retail Contexts}
\label{tab:forecast_compare}
\resizebox{\textwidth}{!}{%
\begin{tabular}{@{}lcccccc@{}}
\toprule
\textbf{Model} & \textbf{Seasonality} & \textbf{Non-Linearity} & \textbf{Multivariate} & \textbf{Data Need} & \textbf{Typical MAPE} & \textbf{Interpretability} \\
\midrule
SARIMA          & Yes (explicit) & No   & No      & Low    & 10--16\% & High \\
Prophet         & Yes (Fourier)  & Weak & Partial & Low    & 9--14\%  & High \\
Random Forest   & Via features   & Yes  & Yes     & Medium & 8--12\%  & Medium \\
XGBoost         & Via features   & Yes  & Yes     & Medium & 7--10\%  & Medium \\
LSTM            & Implicit       & Yes  & Yes     & High   & 6--11\%  & Low \\
\bottomrule
\multicolumn{7}{l}{\footnotesize Sources: \cite{MethodsSalesForecasting2022}; \cite{reportKuijt2022}; \cite{mathematics2023RF}; \cite{1pdf_deep_learning}.}
\end{tabular}}
\end{table}

\subsection{The Customer--Demand Nexus: Integrating Segmentation into Forecasting}

A growing strand of the literature argues that customer-level behavioural features — and in particular CLV segment membership — constitute high-signal exogenous regressors for aggregate demand forecasting. \cite{preprints2026} formalise this argument through a multi-level hierarchical model in which store-level weekly forecasts are conditioned on segment-level purchasing propensity scores derived from the BG/NBD model, achieving a 14\% MAPE reduction relative to unconditional SARIMA benchmarks. The intuition is straightforward: knowing that a store's ``Champions'' segment has unusually high predicted transaction probability in the coming two weeks provides a structurally informative demand signal absent from historical sales lags alone.

\cite{AITU2022} further demonstrate that in fashion retail contexts characterised by strong recency effects — consistent with the PUMA Algeria dataset's Friday peak / Sunday trough pattern — RFM segment transitions within a week provide leading indicators of intra-week demand shifts with a one-to-two day lead time.

\subsection{Gap 3: Identified Literature Gap in Demand Forecasting}

\begin{tcolorbox}[colback=gray!10, colframe=black, title=\textbf{Literature Gap — Demand Forecasting}, fonttitle=\bfseries]
No study in the reviewed literature has evaluated the performance gain from feeding BG/NBD-derived segment-level purchasing propensity scores into a SARIMA or XGBoost forecasting model for a fashion-sportswear distributor in a North African emerging market. The day-of-week demand anomaly documented in the PUMA Algeria dataset — a Sunday crash following Friday and Saturday peaks that diverges from the Western weekend-shopping pattern — constitutes a contextual novelty not addressed in any existing benchmark. This thesis provides the first empirical evaluation of this integrated pipeline on such a dataset.
\end{tcolorbox}

% =============================================================================
\section{Synthesis and Positioning of the Present Thesis}
% =============================================================================

Table~\ref{tab:lit_gap_summary} consolidates the three identified gaps and maps them to the corresponding thesis contributions.

\begin{table}[H]
\centering
\caption{Summary of Literature Gaps and Corresponding Thesis Contributions}
\label{tab:lit_gap_summary}
\begin{tabular}{@{}p{3cm}p{5.5cm}p{5.5cm}@{}}
\toprule
\textbf{Domain} & \textbf{Gap in Existing Literature} & \textbf{Thesis Contribution} \\
\midrule
Customer Segmentation & RFM-K-Means benchmarks confined to Western / East Asian multi-brand e-commerce & Application and validation on SARL Apparel PUMA Algeria mono-brand B2B/B2C dataset (180,000+ rows) \\
\addlinespace
CLV Modelling & BG/NBD never tested in exclusive-distribution emerging-market with mixed channel structure & 90-day holdout validation on PUMA Algeria; assumption diagnostic tests \\
\addlinespace
Demand Forecasting & No CLV-infused forecasting benchmark for North African fashion retail with documented day-of-week anomaly & SARIMA vs.\ XGBoost vs.\ Random Forest vs.\ Prophet multi-model race with segment propensity as exogenous regressor \\
\bottomrule
\end{tabular}
\end{table}

The theoretical contribution of this thesis is the construction and empirical validation of an \textit{integrated customer intelligence pipeline} — segmentation feeding CLV modelling feeding demand forecasting — tailored to the institutional and behavioural specificities of a North African exclusive-brand distributor. The methodological contribution is the systematic evaluation of this pipeline using rigorous holdout protocols, cluster validity indices, and multi-metric forecast benchmarking. Together, these contributions advance the application of AI-driven retail analytics beyond the geographical and institutional boundaries of the existing literature.

% ─────────────────────────────────────────────────────────────────────────────
\section*{Conclusion}
\addcontentsline{toc}{section}{Conclusion}
% ─────────────────────────────────────────────────────────────────────────────

This chapter has reviewed the theoretical and empirical literature across three interconnected domains — customer segmentation, CLV modelling, and demand forecasting — and identified three specific gaps that the present thesis addresses. The RFM framework, while foundational, lacks the optimality guarantees and distributional flexibility of machine-learning clustering. The BG/NBD + Gamma-Gamma probabilistic CLV pipeline represents the current state of the art for non-contractual retail settings but has not been evaluated in North African emerging-market contexts. Demand forecasting has advanced substantially through gradient-boosted tree ensembles, yet the integration of CLV-derived behavioural signals as exogenous regressors remains underexplored. Chapter 3 presents the methodology and empirical results of the integrated pipeline designed to address all three gaps simultaneously.